%%%%%%%%%%%%%%%%%%%%%%%%%%%%%%%%%%%%%%%%%%%%%%%%%%%%%%%%%%%%
% 6.10 RAMSEY THEORY
% The Composition of Advanced Mathematics
%%%%%%%%%%%%%%%%%%%%%%%%%%%%%%%%%%%%%%%%%%%%%%%%%%%%%%%%%%%%

\section{Ramsey Theory}
\label{sec:ramsey-theory}

\prerequisite{
Basic counting, graph theory, complete graphs, graph coloring,
combinations, the pigeonhole principle, induction, and elementary
probability.
}

\learningobjectives{
By the end of this section, the reader should be able to:
\begin{itemize}
    \item explain the central idea of Ramsey theory;
    \item interpret Ramsey problems using edge-colored complete graphs;
    \item define two-color Ramsey numbers;
    \item work with symmetric and asymmetric Ramsey numbers;
    \item prove simple Ramsey bounds using the pigeonhole principle;
    \item prove that \(R(3,3)=6\);
    \item use the standard recursive upper bound for Ramsey numbers;
    \item understand the relationship between Ramsey numbers and
          independent sets;
    \item recognize multicolor Ramsey numbers;
    \item understand the infinite Ramsey theorem;
    \item recognize hypergraph Ramsey theory;
    \item use probabilistic methods to obtain Ramsey lower bounds;
    \item distinguish constructive upper bounds from probabilistic
          lower bounds;
    \item understand Ramsey-type results for arithmetic progressions
          and ordered structures;
    \item recognize the general principle that sufficiently large
          structures must contain ordered substructure.
\end{itemize}
}

%%%%%%%%%%%%%%%%%%%%%%%%%%%%%%%%%%%%%%%%%%%%%%%%%%%%%%%%%%%%
\subsection{What Is Ramsey Theory?}
\label{subsec:what-is-ramsey-theory}
%%%%%%%%%%%%%%%%%%%%%%%%%%%%%%%%%%%%%%%%%%%%%%%%%%%%%%%%%%%%

Ramsey theory studies the principle that sufficiently large discrete
structures must contain highly organized substructures.

Its central philosophy is often summarized as

\[
\boxed{
\text{sufficiently large disorder forces some order}.
}
\]

A typical Ramsey problem asks:

\[
\boxed{
\text{How large must a structure be before a desired configuration is
unavoidable?}
}
\]

This is closely related to extremal combinatorics.

Extremal combinatorics asks how large a structure may become while
avoiding a configuration.

Ramsey theory asks when avoidance becomes impossible.

%%%%%%%%%%%%%%%%%%%%%%%%%%%%%%%%%%%%%%%%%%%%%%%%%%%%%%%%%%%%
\subsection{The Party Problem}
\label{subsec:party-problem}
%%%%%%%%%%%%%%%%%%%%%%%%%%%%%%%%%%%%%%%%%%%%%%%%%%%%%%%%%%%%

A classical example asks:

\begin{quote}
How many people are needed to guarantee that either three people all
know one another or three people are pairwise strangers?
\end{quote}

Represent each person by a vertex.

Color the edge joining two people:

\[
\text{red}
\]

if they know one another, and

\[
\text{blue}
\]

if they do not.

The problem becomes:

\[
\boxed{
\text{How large must a complete graph be to force a monochromatic
triangle?}
}
\]

The answer is \(6\).

%%%%%%%%%%%%%%%%%%%%%%%%%%%%%%%%%%%%%%%%%%%%%%%%%%%%%%%%%%%%
\subsection{Two-Color Ramsey Numbers}
\label{subsec:two-color-ramsey-numbers}
%%%%%%%%%%%%%%%%%%%%%%%%%%%%%%%%%%%%%%%%%%%%%%%%%%%%%%%%%%%%

\begin{definitionbox}{Ramsey Number}
The Ramsey number

\[
\boxed{
R(s,t)
}
\]

is the smallest positive integer \(n\) such that every red-blue coloring
of the edges of \(K_n\) contains either:

\[
\text{a red }K_s,
\]

or

\[
\text{a blue }K_t.
\]
\end{definitionbox}

The notation

\[
R(s,t)
\]

is read ``Ramsey \(s,t\).''

The diagonal Ramsey number is

\[
\boxed{
R(k,k).
}
\]

%%%%%%%%%%%%%%%%%%%%%%%%%%%%%%%%%%%%%%%%%%%%%%%%%%%%%%%%%%%%
\subsection{Symmetry of Ramsey Numbers}
\label{subsec:ramsey-symmetry}
%%%%%%%%%%%%%%%%%%%%%%%%%%%%%%%%%%%%%%%%%%%%%%%%%%%%%%%%%%%%

Interchanging the names red and blue does not change the problem.

Therefore,

\[
\boxed{
R(s,t)=R(t,s).
}
\]

This symmetry reduces the number of cases one must study.

%%%%%%%%%%%%%%%%%%%%%%%%%%%%%%%%%%%%%%%%%%%%%%%%%%%%%%%%%%%%
\subsection{Trivial Ramsey Numbers}
\label{subsec:trivial-ramsey-numbers}
%%%%%%%%%%%%%%%%%%%%%%%%%%%%%%%%%%%%%%%%%%%%%%%%%%%%%%%%%%%%

Several Ramsey numbers follow directly from the definition.

For \(t\geq1\),

\[
\boxed{
R(1,t)=1.
}
\]

For \(t\geq2\),

\[
\boxed{
R(2,t)=t.
}
\]

To see the second identity, any \(K_t\) either contains a red edge or
all edges are blue.

%%%%%%%%%%%%%%%%%%%%%%%%%%%%%%%%%%%%%%%%%%%%%%%%%%%%%%%%%%%%
\subsection{The Ramsey Number \(R(3,3)\)}
\label{subsec:ramsey-33}
%%%%%%%%%%%%%%%%%%%%%%%%%%%%%%%%%%%%%%%%%%%%%%%%%%%%%%%%%%%%

The classical result is

\[
\boxed{
R(3,3)=6.
}
\]

To prove this equality, two things are required:

\begin{enumerate}
    \item show every two-coloring of \(K_6\) contains a monochromatic
          triangle;
    \item construct a coloring of \(K_5\) with no monochromatic
          triangle.
\end{enumerate}

The first gives

\[
R(3,3)\leq6,
\]

while the second gives

\[
R(3,3)>5.
\]

%%%%%%%%%%%%%%%%%%%%%%%%%%%%%%%%%%%%%%%%%%%%%%%%%%%%%%%%%%%%
\subsection{Proof that \(R(3,3)\leq6\)}
\label{subsec:ramsey-33-upper}
%%%%%%%%%%%%%%%%%%%%%%%%%%%%%%%%%%%%%%%%%%%%%%%%%%%%%%%%%%%%

Choose any vertex \(v\) in \(K_6\).

The vertex \(v\) is incident with \(5\) edges.

By the pigeonhole principle, at least

\[
\left\lceil\frac52\right\rceil=3
\]

of these edges have the same color.

Without loss of generality, suppose

\[
va,\quad vb,\quad vc
\]

are red.

Now consider the three edges among

\[
a,b,c.
\]

If any one of them is red, say

\[
ab,
\]

then

\[
v,a,b
\]

form a red triangle.

If none is red, then all three are blue.

Thus

\[
a,b,c
\]

form a blue triangle.

Therefore every red-blue coloring of \(K_6\) contains a monochromatic
triangle.

Hence,

\[
\boxed{
R(3,3)\leq6.
}
\]

%%%%%%%%%%%%%%%%%%%%%%%%%%%%%%%%%%%%%%%%%%%%%%%%%%%%%%%%%%%%
\subsection{Proof that \(R(3,3)>5\)}
\label{subsec:ramsey-33-lower}
%%%%%%%%%%%%%%%%%%%%%%%%%%%%%%%%%%%%%%%%%%%%%%%%%%%%%%%%%%%%

Arrange five vertices around a cycle.

Color the five cycle edges red.

Color the remaining five edges blue.

The red graph is

\[
C_5,
\]

which contains no triangle.

The blue graph is also a \(5\)-cycle, so it also contains no triangle.

Thus there exists a red-blue coloring of \(K_5\) containing no
monochromatic triangle.

Therefore,

\[
\boxed{
R(3,3)>5.
}
\]

Combining both inequalities gives

\[
\boxed{
R(3,3)=6.
}
\]

%%%%%%%%%%%%%%%%%%%%%%%%%%%%%%%%%%%%%%%%%%%%%%%%%%%%%%%%%%%%
\subsection{Worked Example: The Party Problem}
\label{subsec:worked-party-problem}
%%%%%%%%%%%%%%%%%%%%%%%%%%%%%%%%%%%%%%%%%%%%%%%%%%%%%%%%%%%%

\begin{workedexamplebox}{Six People}

Suppose six people attend a gathering.

For every pair, either they know one another or they do not.

Color the corresponding edge red for acquaintance and blue for
strangers.

Since

\[
R(3,3)=6,
\]

the coloring must contain either:

\[
\text{a red triangle},
\]

or

\[
\text{a blue triangle}.
\]

\begin{answerbox}
\[
\boxed{
\text{Among any six people, there are either three mutual acquaintances
or three mutual strangers.}
}
\]
\end{answerbox}

\end{workedexamplebox}

%%%%%%%%%%%%%%%%%%%%%%%%%%%%%%%%%%%%%%%%%%%%%%%%%%%%%%%%%%%%
\subsection{The Fundamental Ramsey Recurrence}
\label{subsec:ramsey-recurrence}
%%%%%%%%%%%%%%%%%%%%%%%%%%%%%%%%%%%%%%%%%%%%%%%%%%%%%%%%%%%%

Ramsey numbers satisfy a useful recursive upper bound.

\begin{theorem}[Ramsey Recurrence]
For integers \(s,t\geq2\),

\[
\boxed{
R(s,t)
\leq
R(s-1,t)+R(s,t-1).
}
\]
\end{theorem}

This recurrence is one of the most basic tools for proving finite Ramsey
numbers exist.

%%%%%%%%%%%%%%%%%%%%%%%%%%%%%%%%%%%%%%%%%%%%%%%%%%%%%%%%%%%%
\subsection{Proof of the Ramsey Recurrence}
\label{subsec:proof-ramsey-recurrence}
%%%%%%%%%%%%%%%%%%%%%%%%%%%%%%%%%%%%%%%%%%%%%%%%%%%%%%%%%%%%

Let

\[
n
=
R(s-1,t)+R(s,t-1).
\]

Consider any red-blue coloring of \(K_n\).

Choose a vertex \(v\).

Among the remaining \(n-1\) vertices, let:

\[
A
=
\text{red neighbors of }v,
\]

and

\[
B
=
\text{blue neighbors of }v.
\]

Then

\[
|A|+|B|=n-1.
\]

Suppose both

\[
|A|<R(s-1,t)
\]

and

\[
|B|<R(s,t-1).
\]

Then

\[
|A|+|B|
\leq
R(s-1,t)-1
+
R(s,t-1)-1
=
n-2,
\]

contradicting

\[
|A|+|B|=n-1.
\]

Therefore at least one of the following holds:

\[
|A|\geq R(s-1,t),
\]

or

\[
|B|\geq R(s,t-1).
\]

If

\[
|A|\geq R(s-1,t),
\]

then the coloring induced on \(A\) contains either a red \(K_{s-1}\) or
a blue \(K_t\).

A red \(K_{s-1}\), together with \(v\), forms a red \(K_s\).

If

\[
|B|\geq R(s,t-1),
\]

the analogous argument produces either a red \(K_s\) or a blue \(K_t\).

Thus,

\[
\boxed{
R(s,t)
\leq
R(s-1,t)+R(s,t-1).
}
\]

%%%%%%%%%%%%%%%%%%%%%%%%%%%%%%%%%%%%%%%%%%%%%%%%%%%%%%%%%%%%
\subsection{Improved Recurrence When Both Terms Are Even}
\label{subsec:ramsey-improved-recurrence}
%%%%%%%%%%%%%%%%%%%%%%%%%%%%%%%%%%%%%%%%%%%%%%%%%%%%%%%%%%%%

A slight refinement states that if both

\[
R(s-1,t)
\]

and

\[
R(s,t-1)
\]

are even, then

\[
\boxed{
R(s,t)
\leq
R(s-1,t)+R(s,t-1)-1.
}
\]

The proof uses parity considerations applied to the red and blue degree
sequences.

This refinement can improve some small Ramsey bounds.

%%%%%%%%%%%%%%%%%%%%%%%%%%%%%%%%%%%%%%%%%%%%%%%%%%%%%%%%%%%%
\subsection{Worked Example: Bounding \(R(3,4)\)}
\label{subsec:worked-r34}
%%%%%%%%%%%%%%%%%%%%%%%%%%%%%%%%%%%%%%%%%%%%%%%%%%%%%%%%%%%%

\begin{workedexamplebox}{Recursive Ramsey Bound}

Using

\[
R(2,4)=4
\]

and

\[
R(3,3)=6,
\]

the recurrence gives

\[
R(3,4)
\leq
R(2,4)+R(3,3).
\]

Therefore,

\[
R(3,4)\leq10.
\]

Because both \(4\) and \(6\) are even, the refined recurrence gives

\[
R(3,4)\leq9.
\]

In fact,

\[
R(3,4)=9.
\]

\begin{answerbox}
\[
\boxed{
R(3,4)=9.
}
\]
\end{answerbox}

\end{workedexamplebox}

%%%%%%%%%%%%%%%%%%%%%%%%%%%%%%%%%%%%%%%%%%%%%%%%%%%%%%%%%%%%
\subsection{Ramsey Numbers and Graph Complements}
\label{subsec:ramsey-graph-complements}
%%%%%%%%%%%%%%%%%%%%%%%%%%%%%%%%%%%%%%%%%%%%%%%%%%%%%%%%%%%%

Ramsey numbers can be expressed without edge colors.

A graph \(G\) on \(n\) vertices and its complement \(\overline{G}\)
partition all possible edges of \(K_n\).

Thus,

\[
R(s,t)
\]

is the smallest \(n\) such that every graph \(G\) on \(n\) vertices
contains either:

\[
\boxed{
K_s\subseteq G
}
\]

or

\[
\boxed{
K_t\subseteq\overline{G}.
}
\]

Since a clique in \(\overline{G}\) corresponds to an independent set in
\(G\), we obtain another interpretation.

%%%%%%%%%%%%%%%%%%%%%%%%%%%%%%%%%%%%%%%%%%%%%%%%%%%%%%%%%%%%
\subsection{Cliques and Independent Sets}
\label{subsec:ramsey-cliques-independent}
%%%%%%%%%%%%%%%%%%%%%%%%%%%%%%%%%%%%%%%%%%%%%%%%%%%%%%%%%%%%

Let

\[
\omega(G)
\]

denote the clique number of \(G\), and let

\[
\alpha(G)
\]

denote its independence number.

Then \(R(s,t)\) is the smallest \(n\) such that every \(n\)-vertex graph
satisfies

\[
\boxed{
\omega(G)\geq s
\quad\text{or}\quad
\alpha(G)\geq t.
}
\]

Thus Ramsey theory guarantees large homogeneous vertex sets.

%%%%%%%%%%%%%%%%%%%%%%%%%%%%%%%%%%%%%%%%%%%%%%%%%%%%%%%%%%%%
\subsection{Diagonal Ramsey Numbers}
\label{subsec:diagonal-ramsey-numbers}
%%%%%%%%%%%%%%%%%%%%%%%%%%%%%%%%%%%%%%%%%%%%%%%%%%%%%%%%%%%%

The diagonal Ramsey number

\[
R(k,k)
\]

is usually abbreviated

\[
\boxed{
R(k).
}
\]

It is the smallest \(n\) such that every red-blue coloring of \(K_n\)
contains a monochromatic \(K_k\).

Equivalently, every graph on \(n\) vertices contains either a clique of
size \(k\) or an independent set of size \(k\).

%%%%%%%%%%%%%%%%%%%%%%%%%%%%%%%%%%%%%%%%%%%%%%%%%%%%%%%%%%%%
\subsection{An Elementary Upper Bound}
\label{subsec:elementary-ramsey-upper}
%%%%%%%%%%%%%%%%%%%%%%%%%%%%%%%%%%%%%%%%%%%%%%%%%%%%%%%%%%%%

Using the Ramsey recurrence repeatedly yields

\[
\boxed{
R(s,t)
\leq
\binom{s+t-2}{s-1}.
}
\]

By symmetry,

\[
\binom{s+t-2}{s-1}
=
\binom{s+t-2}{t-1}.
\]

For the diagonal case,

\[
\boxed{
R(k,k)
\leq
\binom{2k-2}{k-1}.
}
\]

%%%%%%%%%%%%%%%%%%%%%%%%%%%%%%%%%%%%%%%%%%%%%%%%%%%%%%%%%%%%
\subsection{Proof of the Binomial Upper Bound}
\label{subsec:proof-binomial-ramsey-upper}
%%%%%%%%%%%%%%%%%%%%%%%%%%%%%%%%%%%%%%%%%%%%%%%%%%%%%%%%%%%%

We prove by induction that

\[
R(s,t)
\leq
\binom{s+t-2}{s-1}.
\]

The boundary cases are

\[
R(1,t)=1
\]

and

\[
R(s,1)=1.
\]

Assume the result for smaller parameters.

Using the recurrence,

\[
R(s,t)
\leq
R(s-1,t)+R(s,t-1).
\]

Then

\[
R(s,t)
\leq
\binom{s+t-3}{s-2}
+
\binom{s+t-3}{s-1}.
\]

By Pascal's identity,

\[
\boxed{
R(s,t)
\leq
\binom{s+t-2}{s-1}.
}
\]

%%%%%%%%%%%%%%%%%%%%%%%%%%%%%%%%%%%%%%%%%%%%%%%%%%%%%%%%%%%%
\subsection{Exponential Growth}
\label{subsec:ramsey-exponential-growth}
%%%%%%%%%%%%%%%%%%%%%%%%%%%%%%%%%%%%%%%%%%%%%%%%%%%%%%%%%%%%

For the diagonal case,

\[
R(k,k)
\leq
\binom{2k-2}{k-1}.
\]

Since central binomial coefficients grow roughly like

\[
\frac{4^k}{\sqrt{k}},
\]

this shows that diagonal Ramsey numbers grow at most exponentially.

Ramsey theory also has exponential lower bounds.

Thus the correct scale is exponential in \(k\), although the exact
growth remains difficult to determine.

%%%%%%%%%%%%%%%%%%%%%%%%%%%%%%%%%%%%%%%%%%%%%%%%%%%%%%%%%%%%
\subsection{Probabilistic Lower Bounds}
\label{subsec:probabilistic-ramsey-lower}
%%%%%%%%%%%%%%%%%%%%%%%%%%%%%%%%%%%%%%%%%%%%%%%%%%%%%%%%%%%%

Probabilistic combinatorics provides powerful lower bounds.

Color every edge of \(K_n\) independently red or blue with equal
probability.

Fix a set of \(k\) vertices.

It contains

\[
\binom{k}{2}
\]

edges.

The probability that all are red is

\[
2^{-\binom{k}{2}}.
\]

The probability that all are blue is the same.

Therefore,

\[
\boxed{
\mathbb{P}
(\text{fixed }K_k\text{ monochromatic})
=
2^{1-\binom{k}{2}}.
}
\]

%%%%%%%%%%%%%%%%%%%%%%%%%%%%%%%%%%%%%%%%%%%%%%%%%%%%%%%%%%%%
\subsection{First-Moment Ramsey Bound}
\label{subsec:first-moment-ramsey}
%%%%%%%%%%%%%%%%%%%%%%%%%%%%%%%%%%%%%%%%%%%%%%%%%%%%%%%%%%%%

There are

\[
\binom nk
\]

possible \(k\)-vertex subsets.

Let \(X\) count monochromatic \(K_k\)'s.

Then

\[
\boxed{
\mathbb{E}[X]
=
\binom nk
2^{1-\binom{k}{2}}.
}
\]

If

\[
\binom nk
2^{1-\binom{k}{2}}
<1,
\]

then some coloring satisfies

\[
X=0.
\]

Therefore,

\[
\boxed{
R(k,k)>n.
}
\]

%%%%%%%%%%%%%%%%%%%%%%%%%%%%%%%%%%%%%%%%%%%%%%%%%%%%%%%%%%%%
\subsection{Worked Example: Probabilistic Lower Bound}
\label{subsec:worked-probabilistic-ramsey-bound}
%%%%%%%%%%%%%%%%%%%%%%%%%%%%%%%%%%%%%%%%%%%%%%%%%%%%%%%%%%%%

\begin{workedexamplebox}{A General Existence Criterion}

Suppose \(n\) and \(k\) satisfy

\[
\binom nk
2^{1-\binom{k}{2}}
<1.
\]

Randomly color the edges of \(K_n\) red or blue.

Let \(X\) be the number of monochromatic copies of \(K_k\).

Then

\[
\mathbb{E}[X]
=
\binom nk
2^{1-\binom{k}{2}}
<1.
\]

Since \(X\) is a nonnegative integer-valued random variable, there must
be a coloring for which

\[
X=0.
\]

\begin{answerbox}
\[
\boxed{
R(k,k)>n.
}
\]
\end{answerbox}

\end{workedexamplebox}

%%%%%%%%%%%%%%%%%%%%%%%%%%%%%%%%%%%%%%%%%%%%%%%%%%%%%%%%%%%%
\subsection{Classical Asymptotic Lower Bound}
\label{subsec:classical-ramsey-lower}
%%%%%%%%%%%%%%%%%%%%%%%%%%%%%%%%%%%%%%%%%%%%%%%%%%%%%%%%%%%%

The probabilistic method yields a bound of the form

\[
\boxed{
R(k,k)
>
c\,k\,2^{k/2}
}
\]

for an absolute positive constant \(c\).

The important feature is the exponential term

\[
2^{k/2}.
\]

Thus diagonal Ramsey numbers grow at least exponentially.

%%%%%%%%%%%%%%%%%%%%%%%%%%%%%%%%%%%%%%%%%%%%%%%%%%%%%%%%%%%%
\subsection{The Gap Between Upper and Lower Bounds}
\label{subsec:ramsey-gap}
%%%%%%%%%%%%%%%%%%%%%%%%%%%%%%%%%%%%%%%%%%%%%%%%%%%%%%%%%%%%

Elementary arguments give roughly

\[
R(k,k)
\lesssim
4^k,
\]

while probabilistic methods give roughly

\[
R(k,k)
\gtrsim
2^{k/2}.
\]

The precise asymptotic growth of

\[
R(k,k)
\]

remains a central problem.

Even modest improvements to these bounds can be mathematically
significant.

%%%%%%%%%%%%%%%%%%%%%%%%%%%%%%%%%%%%%%%%%%%%%%%%%%%%%%%%%%%%
\subsection{Off-Diagonal Ramsey Numbers}
\label{subsec:off-diagonal-ramsey}
%%%%%%%%%%%%%%%%%%%%%%%%%%%%%%%%%%%%%%%%%%%%%%%%%%%%%%%%%%%%

When

\[
s\neq t,
\]

the Ramsey number

\[
R(s,t)
\]

is called off-diagonal.

These numbers can behave differently from diagonal Ramsey numbers.

For example,

\[
R(3,t)
\]

has been studied extensively because triangle avoidance gives strong
connections with extremal graph theory and random graphs.

%%%%%%%%%%%%%%%%%%%%%%%%%%%%%%%%%%%%%%%%%%%%%%%%%%%%%%%%%%%%
\subsection{Multicolor Ramsey Numbers}
\label{subsec:multicolor-ramsey}
%%%%%%%%%%%%%%%%%%%%%%%%%%%%%%%%%%%%%%%%%%%%%%%%%%%%%%%%%%%%

Ramsey theory generalizes to more than two colors.

\begin{definitionbox}{Multicolor Ramsey Number}
The number

\[
\boxed{
R(k_1,k_2,\ldots,k_r)
}
\]

is the smallest \(n\) such that every coloring of the edges of \(K_n\)
with \(r\) colors contains a monochromatic \(K_{k_i}\) in color \(i\)
for some \(i\).
\end{definitionbox}

When

\[
k_1=\cdots=k_r=k,
\]

one often writes

\[
\boxed{
R_r(k).
}
\]

%%%%%%%%%%%%%%%%%%%%%%%%%%%%%%%%%%%%%%%%%%%%%%%%%%%%%%%%%%%%
\subsection{Three-Color Example}
\label{subsec:three-color-ramsey}
%%%%%%%%%%%%%%%%%%%%%%%%%%%%%%%%%%%%%%%%%%%%%%%%%%%%%%%%%%%%

Suppose the edges of a complete graph are colored:

\[
\text{red},
\qquad
\text{blue},
\qquad
\text{green}.
\]

The number

\[
R(3,3,3)
\]

is the least \(n\) forcing a monochromatic triangle in one of the three
colors.

The principle remains unchanged:

\[
\boxed{
\text{enough vertices force monochromatic structure}.
}
\]

%%%%%%%%%%%%%%%%%%%%%%%%%%%%%%%%%%%%%%%%%%%%%%%%%%%%%%%%%%%%
\subsection{Graph Ramsey Numbers}
\label{subsec:graph-ramsey-numbers}
%%%%%%%%%%%%%%%%%%%%%%%%%%%%%%%%%%%%%%%%%%%%%%%%%%%%%%%%%%%%

The complete graphs in the definition may be replaced by arbitrary
target graphs.

For graphs \(G\) and \(H\), define

\[
\boxed{
R(G,H)
}
\]

to be the least \(n\) such that every red-blue coloring of \(K_n\)
contains either:

\[
\text{a red copy of }G,
\]

or

\[
\text{a blue copy of }H.
\]

Classical Ramsey numbers satisfy

\[
\boxed{
R(s,t)=R(K_s,K_t).
}
\]

%%%%%%%%%%%%%%%%%%%%%%%%%%%%%%%%%%%%%%%%%%%%%%%%%%%%%%%%%%%%
\subsection{Infinite Ramsey Theory}
\label{subsec:infinite-ramsey}
%%%%%%%%%%%%%%%%%%%%%%%%%%%%%%%%%%%%%%%%%%%%%%%%%%%%%%%%%%%%

Ramsey theory also has an infinite form.

\begin{theorem}[Infinite Ramsey Theorem]
Let the \(k\)-element subsets of an infinite set be colored using
finitely many colors.

Then there exists an infinite subset whose \(k\)-element subsets all
have the same color.
\end{theorem}

For \(k=2\), every finite coloring of the edges of an infinite complete
graph contains an infinite monochromatic complete subgraph.

%%%%%%%%%%%%%%%%%%%%%%%%%%%%%%%%%%%%%%%%%%%%%%%%%%%%%%%%%%%%
\subsection{Why the Infinite Theorem Is Strong}
\label{subsec:why-infinite-ramsey-strong}
%%%%%%%%%%%%%%%%%%%%%%%%%%%%%%%%%%%%%%%%%%%%%%%%%%%%%%%%%%%%

A finite Ramsey theorem guarantees a finite structured subset inside a
sufficiently large finite structure.

The infinite theorem guarantees an infinite structured subset.

Thus arbitrary finite coloring cannot destroy all large-scale
homogeneity.

This theorem has important consequences in:

\begin{itemize}
    \item combinatorics;
    \item logic;
    \item set theory;
    \item topology.
\end{itemize}

%%%%%%%%%%%%%%%%%%%%%%%%%%%%%%%%%%%%%%%%%%%%%%%%%%%%%%%%%%%%
\subsection{Hypergraph Ramsey Theory}
\label{subsec:hypergraph-ramsey}
%%%%%%%%%%%%%%%%%%%%%%%%%%%%%%%%%%%%%%%%%%%%%%%%%%%%%%%%%%%%

Ramsey theory extends from pairs to \(r\)-element subsets.

Color every \(r\)-element subset of an \(n\)-element set.

A hypergraph Ramsey number asks how large \(n\) must be before a large
subset has all of its \(r\)-element subsets the same color.

Hypergraph Ramsey numbers grow much faster than ordinary graph Ramsey
numbers.

%%%%%%%%%%%%%%%%%%%%%%%%%%%%%%%%%%%%%%%%%%%%%%%%%%%%%%%%%%%%
\subsection{Uniform Hypergraph Ramsey Numbers}
\label{subsec:uniform-hypergraph-ramsey}
%%%%%%%%%%%%%%%%%%%%%%%%%%%%%%%%%%%%%%%%%%%%%%%%%%%%%%%%%%%%

For integers

\[
r,k,q,
\]

one may define an \(r\)-uniform \(q\)-color Ramsey number.

The objects being colored are

\[
\binom{[n]}{r}.
\]

The goal is to force a \(k\)-element subset \(S\) such that all members
of

\[
\binom{S}{r}
\]

receive the same color.

The graph case corresponds to

\[
r=2.
\]

%%%%%%%%%%%%%%%%%%%%%%%%%%%%%%%%%%%%%%%%%%%%%%%%%%%%%%%%%%%%
\subsection{Ramsey Theory for Sequences}
\label{subsec:ramsey-sequences}
%%%%%%%%%%%%%%%%%%%%%%%%%%%%%%%%%%%%%%%%%%%%%%%%%%%%%%%%%%%%

Ramsey-type phenomena occur outside graphs.

A famous example is the Erd\H{o}s--Szekeres theorem.

\begin{theorem}[Erd\H{o}s--Szekeres Monotone Subsequence Theorem]
Every sequence of

\[
(r-1)(s-1)+1
\]

distinct real numbers contains either:

\[
\text{an increasing subsequence of length }r,
\]

or

\[
\text{a decreasing subsequence of length }s.
\]
\end{theorem}

This is another example of unavoidable order.

%%%%%%%%%%%%%%%%%%%%%%%%%%%%%%%%%%%%%%%%%%%%%%%%%%%%%%%%%%%%
\subsection{Worked Example: Monotone Subsequences}
\label{subsec:worked-erdos-szekeres}
%%%%%%%%%%%%%%%%%%%%%%%%%%%%%%%%%%%%%%%%%%%%%%%%%%%%%%%%%%%%

\begin{workedexamplebox}{Increasing or Decreasing Subsequence}

Take a sequence of

\[
10
\]

distinct real numbers.

Choose

\[
r=s=4.
\]

Since

\[
(r-1)(s-1)+1
=
3\cdot3+1
=
10,
\]

the Erd\H{o}s--Szekeres theorem applies.

\begin{answerbox}
\[
\boxed{
\text{Every sequence of }10\text{ distinct numbers contains an
increasing or decreasing subsequence of length }4.
}
\]
\end{answerbox}

\end{workedexamplebox}

%%%%%%%%%%%%%%%%%%%%%%%%%%%%%%%%%%%%%%%%%%%%%%%%%%%%%%%%%%%%
\subsection{Geometric Ramsey Phenomena}
\label{subsec:geometric-ramsey}
%%%%%%%%%%%%%%%%%%%%%%%%%%%%%%%%%%%%%%%%%%%%%%%%%%%%%%%%%%%%

The Erd\H{o}s--Szekeres theorem also appears in geometry.

One classical form states that sufficiently many points in the plane in
general position contain a subset forming the vertices of a convex
polygon of prescribed size.

This illustrates the general Ramsey principle:

\[
\boxed{
\text{large geometric configuration}
\Longrightarrow
\text{ordered geometric subconfiguration}.
}
\]

%%%%%%%%%%%%%%%%%%%%%%%%%%%%%%%%%%%%%%%%%%%%%%%%%%%%%%%%%%%%
\subsection{Van der Waerden's Theorem}
\label{subsec:van-der-waerden}
%%%%%%%%%%%%%%%%%%%%%%%%%%%%%%%%%%%%%%%%%%%%%%%%%%%%%%%%%%%%

Ramsey phenomena also occur in arithmetic.

\begin{theorem}[Van der Waerden's Theorem]
For every positive integers \(r\) and \(k\), there exists an integer
\(N\) such that every coloring of

\[
\{1,2,\ldots,N\}
\]

with \(r\) colors contains a monochromatic arithmetic progression of
length \(k\).
\end{theorem}

An arithmetic progression has the form

\[
\boxed{
a,\,
a+d,\,
a+2d,\,
\ldots,\,
a+(k-1)d.
}
\]

%%%%%%%%%%%%%%%%%%%%%%%%%%%%%%%%%%%%%%%%%%%%%%%%%%%%%%%%%%%%
\subsection{Van der Waerden Numbers}
\label{subsec:van-der-waerden-numbers}
%%%%%%%%%%%%%%%%%%%%%%%%%%%%%%%%%%%%%%%%%%%%%%%%%%%%%%%%%%%%

The smallest such \(N\) is called a van der Waerden number.

One notation is

\[
\boxed{
W(r,k).
}
\]

Thus,

\[
W(r,k)
\]

is the minimum \(N\) such that every \(r\)-coloring of

\[
[N]
\]

contains a monochromatic \(k\)-term arithmetic progression.

%%%%%%%%%%%%%%%%%%%%%%%%%%%%%%%%%%%%%%%%%%%%%%%%%%%%%%%%%%%%
\subsection{Schur's Theorem}
\label{subsec:schurs-theorem}
%%%%%%%%%%%%%%%%%%%%%%%%%%%%%%%%%%%%%%%%%%%%%%%%%%%%%%%%%%%%

Another arithmetic Ramsey theorem concerns equations.

\begin{theorem}[Schur's Theorem]
For every finite coloring of the positive integers, there exist
positive integers \(x,y,z\) of the same color satisfying

\[
\boxed{
x+y=z.
}
\]
\end{theorem}

Finite versions give Schur numbers.

This theorem connects Ramsey theory with additive combinatorics.

%%%%%%%%%%%%%%%%%%%%%%%%%%%%%%%%%%%%%%%%%%%%%%%%%%%%%%%%%%%%
\subsection{Ramsey Theory and the Pigeonhole Principle}
\label{subsec:ramsey-pigeonhole}
%%%%%%%%%%%%%%%%%%%%%%%%%%%%%%%%%%%%%%%%%%%%%%%%%%%%%%%%%%%%

The pigeonhole principle can be viewed as a simple Ramsey-type theorem.

If sufficiently many objects are placed into finitely many classes,
some class must contain many objects.

Ramsey theory extends this idea from individual objects to relationships
among objects.

Thus,

\[
\boxed{
\text{pigeonhole principle}
\quad\text{is a precursor to}\quad
\text{Ramsey theory}.
}
\]

%%%%%%%%%%%%%%%%%%%%%%%%%%%%%%%%%%%%%%%%%%%%%%%%%%%%%%%%%%%%
\subsection{Ramsey Theory and Extremal Combinatorics}
\label{subsec:ramsey-extremal-connection}
%%%%%%%%%%%%%%%%%%%%%%%%%%%%%%%%%%%%%%%%%%%%%%%%%%%%%%%%%%%%

Extremal and Ramsey theory ask complementary questions.

Extremal theory asks:

\[
\boxed{
\text{How large can a structure be while avoiding }H?
}
\]

Ramsey theory asks:

\[
\boxed{
\text{How large must the ambient structure be before avoiding all
specified patterns becomes impossible?}
}
\]

Extremal numbers often contribute to Ramsey bounds.

%%%%%%%%%%%%%%%%%%%%%%%%%%%%%%%%%%%%%%%%%%%%%%%%%%%%%%%%%%%%
\subsection{Ramsey Theory and Probability}
\label{subsec:ramsey-probability-connection}
%%%%%%%%%%%%%%%%%%%%%%%%%%%%%%%%%%%%%%%%%%%%%%%%%%%%%%%%%%%%

Upper bounds are often proved deterministically.

Lower bounds are often obtained by constructing colorings or graphs
with no large homogeneous substructure.

Random constructions are exceptionally effective.

This produces one of the clearest examples of the probabilistic method:

\[
\boxed{
\text{random coloring}
\Longrightarrow
\text{positive probability of avoidance}
\Longrightarrow
\text{deterministic coloring exists}.
}
\]

%%%%%%%%%%%%%%%%%%%%%%%%%%%%%%%%%%%%%%%%%%%%%%%%%%%%%%%%%%%%
\subsection{Ramsey Graphs}
\label{subsec:ramsey-graphs}
%%%%%%%%%%%%%%%%%%%%%%%%%%%%%%%%%%%%%%%%%%%%%%%%%%%%%%%%%%%%

A Ramsey graph is informally a graph having neither a large clique nor a
large independent set.

For a graph \(G\),

\[
\omega(G)
\]

measures its largest clique, while

\[
\alpha(G)
\]

measures its largest independent set.

A good Ramsey graph keeps both quantities small relative to the number
of vertices.

Random graphs naturally produce such behavior.

%%%%%%%%%%%%%%%%%%%%%%%%%%%%%%%%%%%%%%%%%%%%%%%%%%%%%%%%%%%%
\subsection{Pseudorandom Graphs and Ramsey Theory}
\label{subsec:pseudorandom-ramsey}
%%%%%%%%%%%%%%%%%%%%%%%%%%%%%%%%%%%%%%%%%%%%%%%%%%%%%%%%%%%%

Explicit Ramsey constructions often attempt to imitate random graphs.

A pseudorandom graph may have:

\begin{itemize}
    \item nearly uniform edge distribution;
    \item small nontrivial eigenvalues;
    \item expected-like subgraph counts;
    \item no unusually large clique or independent set.
\end{itemize}

Thus algebraic and spectral methods can provide explicit substitutes for
probabilistic constructions.

%%%%%%%%%%%%%%%%%%%%%%%%%%%%%%%%%%%%%%%%%%%%%%%%%%%%%%%%%%%%
\subsection{Ramsey Multiplicity}
\label{subsec:ramsey-multiplicity}
%%%%%%%%%%%%%%%%%%%%%%%%%%%%%%%%%%%%%%%%%%%%%%%%%%%%%%%%%%%%

Classical Ramsey theory asks whether at least one monochromatic
configuration is unavoidable.

Ramsey multiplicity asks:

\[
\boxed{
\text{How many monochromatic copies must necessarily occur?}
}
\]

For example, once \(n\) is large, every two-coloring of \(K_n\)
contains many monochromatic triangles.

This strengthens pure existence into quantitative enumeration.

%%%%%%%%%%%%%%%%%%%%%%%%%%%%%%%%%%%%%%%%%%%%%%%%%%%%%%%%%%%%
\subsection{Induced Ramsey Theory}
\label{subsec:induced-ramsey}
%%%%%%%%%%%%%%%%%%%%%%%%%%%%%%%%%%%%%%%%%%%%%%%%%%%%%%%%%%%%

Classical Ramsey theory seeks monochromatic subgraphs.

Induced Ramsey theory additionally requires that the copy preserve both
edges and nonedges.

Thus an induced copy imposes a more restrictive structural condition.

Induced Ramsey numbers are therefore generally harder to analyze.

%%%%%%%%%%%%%%%%%%%%%%%%%%%%%%%%%%%%%%%%%%%%%%%%%%%%%%%%%%%%
\subsection{Ordered Ramsey Theory}
\label{subsec:ordered-ramsey}
%%%%%%%%%%%%%%%%%%%%%%%%%%%%%%%%%%%%%%%%%%%%%%%%%%%%%%%%%%%%

An ordered graph includes a fixed linear order on its vertices.

An ordered Ramsey problem requires the monochromatic copy to preserve
both:

\[
\text{adjacency}
\]

and

\[
\text{vertex order}.
\]

Ordered Ramsey theory connects combinatorics with permutations,
sequences, and geometric configurations.

%%%%%%%%%%%%%%%%%%%%%%%%%%%%%%%%%%%%%%%%%%%%%%%%%%%%%%%%%%%%
\subsection{Canonical Ramsey Theory}
\label{subsec:canonical-ramsey}
%%%%%%%%%%%%%%%%%%%%%%%%%%%%%%%%%%%%%%%%%%%%%%%%%%%%%%%%%%%%

Canonical Ramsey theory replaces the goal of obtaining one color with
the goal of obtaining a simple structural pattern in the coloring.

A large enough colored structure may contain a subset on which the
coloring behaves according to one of a few canonical rules.

Thus even when full monochromaticity is too restrictive, strong
regularity may still be unavoidable.

%%%%%%%%%%%%%%%%%%%%%%%%%%%%%%%%%%%%%%%%%%%%%%%%%%%%%%%%%%%%
\subsection{Common Proof Strategy}
\label{subsec:ramsey-proof-strategy}
%%%%%%%%%%%%%%%%%%%%%%%%%%%%%%%%%%%%%%%%%%%%%%%%%%%%%%%%%%%%

A standard Ramsey proof often follows one of several patterns.

\begin{enumerate}
    \item choose one vertex or element;
    \item partition the remaining structure according to colors or
          relationships;
    \item apply the pigeonhole principle;
    \item invoke a smaller Ramsey result;
    \item extend the smaller monochromatic configuration by the chosen
          vertex.
\end{enumerate}

This is the mechanism behind the recurrence

\[
R(s,t)
\leq
R(s-1,t)+R(s,t-1).
\]

%%%%%%%%%%%%%%%%%%%%%%%%%%%%%%%%%%%%%%%%%%%%%%%%%%%%%%%%%%%%
\subsection{Upper-Bound Strategy}
\label{subsec:ramsey-upper-bound-strategy}
%%%%%%%%%%%%%%%%%%%%%%%%%%%%%%%%%%%%%%%%%%%%%%%%%%%%%%%%%%%%

To prove

\[
R(s,t)\leq N,
\]

one must show:

\[
\boxed{
\text{every red-blue coloring of }K_N
\text{ contains the desired monochromatic configuration}.
}
\]

This is a universal statement.

A single coloring is not enough.

Every possible coloring must be handled.

%%%%%%%%%%%%%%%%%%%%%%%%%%%%%%%%%%%%%%%%%%%%%%%%%%%%%%%%%%%%
\subsection{Lower-Bound Strategy}
\label{subsec:ramsey-lower-bound-strategy}
%%%%%%%%%%%%%%%%%%%%%%%%%%%%%%%%%%%%%%%%%%%%%%%%%%%%%%%%%%%%

To prove

\[
R(s,t)>N,
\]

one only needs to produce one coloring of

\[
K_N
\]

that avoids both:

\[
\text{red }K_s
\]

and

\[
\text{blue }K_t.
\]

Thus lower bounds are existential construction problems.

They may be established by:

\begin{itemize}
    \item explicit construction;
    \item computer search;
    \item algebraic construction;
    \item probabilistic construction.
\end{itemize}

%%%%%%%%%%%%%%%%%%%%%%%%%%%%%%%%%%%%%%%%%%%%%%%%%%%%%%%%%%%%
\subsection{Exact Ramsey Numbers}
\label{subsec:exact-ramsey-numbers}
%%%%%%%%%%%%%%%%%%%%%%%%%%%%%%%%%%%%%%%%%%%%%%%%%%%%%%%%%%%%

Exact Ramsey numbers are notoriously difficult to determine.

Even relatively small parameters may require substantial structural
arguments and computation.

The difficulty arises because one must simultaneously prove:

\[
\boxed{
\text{a universal upper bound}
}
\]

and

\[
\boxed{
\text{a matching construction for the lower bound}.
}
\]

This illustrates a recurring theme in discrete mathematics: finite
problems can still be extremely difficult.

%%%%%%%%%%%%%%%%%%%%%%%%%%%%%%%%%%%%%%%%%%%%%%%%%%%%%%%%%%%%
\subsection{Applications}
\label{subsec:ramsey-theory-applications}
%%%%%%%%%%%%%%%%%%%%%%%%%%%%%%%%%%%%%%%%%%%%%%%%%%%%%%%%%%%%

\begin{applicationbox}

\textbf{Graph Theory.}
Ramsey numbers quantify unavoidable cliques and independent sets.

\medskip

\textbf{Extremal Combinatorics.}
Forbidden-subgraph bounds and Ramsey thresholds are closely linked.

\medskip

\textbf{Probability.}
Random colorings and random graphs provide major Ramsey lower bounds.

\medskip

\textbf{Number Theory.}
Van der Waerden's theorem and Schur's theorem reveal unavoidable
arithmetic patterns.

\medskip

\textbf{Geometry.}
Erd\H{o}s--Szekeres-type results force convex and monotone structures.

\medskip

\textbf{Logic.}
Infinite Ramsey theory has major consequences in combinatorial and
model-theoretic logic.

\medskip

\textbf{Computer Science.}
Ramsey arguments appear in complexity theory, communication complexity,
data structures, and lower-bound proofs.

\medskip

\textbf{Pseudorandomness.}
Explicit Ramsey graphs are closely related to pseudorandom
constructions.

\end{applicationbox}

%%%%%%%%%%%%%%%%%%%%%%%%%%%%%%%%%%%%%%%%%%%%%%%%%%%%%%%%%%%%
\subsection{Common Errors}
\label{subsec:ramsey-common-errors}
%%%%%%%%%%%%%%%%%%%%%%%%%%%%%%%%%%%%%%%%%%%%%%%%%%%%%%%%%%%%

\subsubsection{Confusing a Ramsey Number with a Count of Colorings}

The number

\[
R(s,t)
\]

is a threshold number of vertices.

It does not count the number of colorings.

%%%%%%%%%%%%%%%%%%%%%%%%%%%%%%%%%%%%%%%%%%%%%%%%%%%%%%%%%%%%
\subsubsection{Proving Only the Upper Bound}

To prove an exact identity

\[
R(s,t)=N,
\]

one must prove both:

\[
R(s,t)\leq N
\]

and

\[
R(s,t)>N-1.
\]

%%%%%%%%%%%%%%%%%%%%%%%%%%%%%%%%%%%%%%%%%%%%%%%%%%%%%%%%%%%%
\subsubsection{Using One Coloring to Prove an Upper Bound}

An upper bound must hold for every coloring.

One example coloring can only help prove a lower bound.

%%%%%%%%%%%%%%%%%%%%%%%%%%%%%%%%%%%%%%%%%%%%%%%%%%%%%%%%%%%%
\subsubsection{Forgetting the Complete Ambient Graph}

Classical graph Ramsey numbers color all edges of

\[
K_n.
\]

The underlying ambient graph is complete.

%%%%%%%%%%%%%%%%%%%%%%%%%%%%%%%%%%%%%%%%%%%%%%%%%%%%%%%%%%%%
\subsubsection{Confusing Clique and Independent Set Interpretations}

A blue clique in the two-color formulation corresponds to an independent
set only after one color is interpreted as edges of \(G\) and the other
as edges of \(\overline{G}\).

%%%%%%%%%%%%%%%%%%%%%%%%%%%%%%%%%%%%%%%%%%%%%%%%%%%%%%%%%%%%
\subsubsection{Assuming the Ramsey Recurrence Is an Equality}

In general,

\[
R(s,t)
\leq
R(s-1,t)+R(s,t-1)
\]

is an upper bound.

It is not generally an equality.

%%%%%%%%%%%%%%%%%%%%%%%%%%%%%%%%%%%%%%%%%%%%%%%%%%%%%%%%%%%%
\subsubsection{Assuming Expected Number Below One Means Every Coloring Works}

If

\[
\mathbb{E}[X]<1,
\]

then at least one coloring has \(X=0\).

It does not imply that every coloring has \(X=0\).

%%%%%%%%%%%%%%%%%%%%%%%%%%%%%%%%%%%%%%%%%%%%%%%%%%%%%%%%%%%%
\subsubsection{Confusing Finite and Infinite Ramsey Results}

Finite Ramsey theory provides finite homogeneous subsets once the
ambient size is sufficiently large.

Infinite Ramsey theory guarantees an infinite homogeneous subset in an
infinite structure.

%%%%%%%%%%%%%%%%%%%%%%%%%%%%%%%%%%%%%%%%%%%%%%%%%%%%%%%%%%%%
\subsection{Exercises}
\label{subsec:ramsey-theory-exercises}
%%%%%%%%%%%%%%%%%%%%%%%%%%%%%%%%%%%%%%%%%%%%%%%%%%%%%%%%%%%%

\begin{exercise}[1]
Define the Ramsey number

\[
R(s,t).
\]
\end{exercise}

\begin{exercise}[1]
Explain why

\[
R(s,t)=R(t,s).
\]
\end{exercise}

\begin{exercise}[1]
Show that

\[
R(2,t)=t.
\]
\end{exercise}

\begin{exercise}[1]
State the value of

\[
R(3,3).
\]
\end{exercise}

\begin{exercise}[1]
Explain the party interpretation of

\[
R(3,3)=6.
\]
\end{exercise}

\begin{exercise}[1]
What does a red triangle represent in the acquaintance version of the
party problem?
\end{exercise}

\begin{exercise}[1]
What does a blue triangle represent?
\end{exercise}

\begin{exercise}[1]
Define a diagonal Ramsey number.
\end{exercise}

\begin{exercise}[1]
Define an off-diagonal Ramsey number.
\end{exercise}

\begin{exercise}[1]
State the basic Ramsey recurrence.
\end{exercise}

\begin{exercise}[2]
Prove

\[
R(3,3)\leq6.
\]
\end{exercise}

\begin{exercise}[2]
Construct a red-blue coloring of \(K_5\) containing no monochromatic
triangle.
\end{exercise}

\begin{exercise}[2]
Use the previous two exercises to prove

\[
R(3,3)=6.
\]
\end{exercise}

\begin{exercise}[2]
Use the recurrence to show

\[
R(3,4)
\leq
R(2,4)+R(3,3).
\]
\end{exercise}

\begin{exercise}[2]
Use the basic recurrence to obtain an upper bound for

\[
R(4,4).
\]
\end{exercise}

\begin{exercise}[2]
Interpret \(R(s,t)\) using cliques and independent sets in ordinary
graphs.
\end{exercise}

\begin{exercise}[2]
Show that

\[
R(s,t)
\]

is finite for all positive integers \(s,t\) using induction and the
Ramsey recurrence.
\end{exercise}

\begin{exercise}[2]
Explain why the existence of \(R(s,t)\) is not obvious from the
definition alone.
\end{exercise}

\begin{exercise}[2]
State the meaning of

\[
R(G,H).
\]
\end{exercise}

\begin{exercise}[2]
Define the multicolor Ramsey number

\[
R(k_1,\ldots,k_r).
\]
\end{exercise}

\begin{exercise}[3]
Prove the Ramsey recurrence

\[
R(s,t)
\leq
R(s-1,t)+R(s,t-1).
\]
\end{exercise}

\begin{exercise}[3]
Use induction and Pascal's identity to prove

\[
R(s,t)
\leq
\binom{s+t-2}{s-1}.
\]
\end{exercise}

\begin{exercise}[3]
Deduce that

\[
R(k,k)
\leq
\binom{2k-2}{k-1}.
\]
\end{exercise}

\begin{exercise}[3]
Use the central binomial coefficient estimate to show that the previous
upper bound is exponential in \(k\).
\end{exercise}

\begin{exercise}[3]
In a random red-blue coloring of \(K_n\), compute the probability that a
fixed \(K_k\) is monochromatic.
\end{exercise}

\begin{exercise}[3]
Show that the expected number of monochromatic \(K_k\)'s is

\[
\binom nk
2^{1-\binom{k}{2}}.
\]
\end{exercise}

\begin{exercise}[3]
Prove that if

\[
\binom nk
2^{1-\binom{k}{2}}
<1,
\]

then

\[
R(k,k)>n.
\]
\end{exercise}

\begin{exercise}[3]
Explain why the probabilistic proof above may prove existence without
giving an explicit coloring.
\end{exercise}

\begin{exercise}[3]
For a graph \(G\), express the Ramsey condition using

\[
\omega(G)
\]

and

\[
\alpha(G).
\]
\end{exercise}

\begin{exercise}[3]
Prove the Erd\H{o}s--Szekeres monotone subsequence theorem using the
pigeonhole principle.
\end{exercise}

\begin{exercise}[3]
Apply the Erd\H{o}s--Szekeres theorem with

\[
r=5,
\qquad
s=4
\]

to determine the sequence length that guarantees one of the two
monotone subsequences.
\end{exercise}

\begin{exercise}[3]
Explain how Ramsey theory and extremal combinatorics are complementary.
\end{exercise}

\begin{exercise}[3]
Explain why random graphs naturally produce examples of Ramsey graphs.
\end{exercise}

\begin{exercise}[4]
Study a sharper probabilistic lower bound for

\[
R(k,k).
\]
\end{exercise}

\begin{exercise}[4]
Investigate known techniques for estimating

\[
R(3,t).
\]
\end{exercise}

\begin{exercise}[4]
Study multicolor Ramsey numbers and derive a recursive upper bound.
\end{exercise}

\begin{exercise}[4]
Prove the infinite Ramsey theorem for pairs.
\end{exercise}

\begin{exercise}[4]
Investigate hypergraph Ramsey numbers and compare their growth with
ordinary graph Ramsey numbers.
\end{exercise}

\begin{exercise}[4]
Study Van der Waerden's theorem and compute several small van der
Waerden numbers.
\end{exercise}

\begin{exercise}[4]
Investigate Schur's theorem and its relationship with additive
combinatorics.
\end{exercise}

\begin{exercise}[4]
Study induced Ramsey numbers.
\end{exercise}

\begin{exercise}[4]
Investigate explicit constructions of Ramsey graphs.
\end{exercise}

\begin{exercise}[4]
Study Ramsey multiplicity and determine the minimum number of
monochromatic triangles in a two-coloring of \(K_n\) for small \(n\).
\end{exercise}

%%%%%%%%%%%%%%%%%%%%%%%%%%%%%%%%%%%%%%%%%%%%%%%%%%%%%%%%%%%%
\subsection{Conceptual Summary}
\label{subsec:ramsey-theory-summary}
%%%%%%%%%%%%%%%%%%%%%%%%%%%%%%%%%%%%%%%%%%%%%%%%%%%%%%%%%%%%

Ramsey theory studies unavoidable order in sufficiently large discrete
structures.

The Ramsey number

\[
\boxed{
R(s,t)
}
\]

is the least \(n\) such that every red-blue coloring of \(K_n\) contains
a red \(K_s\) or a blue \(K_t\).

The classical party theorem is

\[
\boxed{
R(3,3)=6.
}
\]

Ramsey numbers satisfy

\[
\boxed{
R(s,t)
\leq
R(s-1,t)+R(s,t-1).
}
\]

This yields the general bound

\[
\boxed{
R(s,t)
\leq
\binom{s+t-2}{s-1}.
}
\]

In ordinary graph language,

\[
R(s,t)
\]

is the smallest \(n\) such that every \(n\)-vertex graph contains either

\[
K_s
\]

or an independent set of size \(t\).

The probabilistic method gives lower bounds.

For a random red-blue coloring,

\[
\boxed{
\mathbb{E}[\#\text{ monochromatic }K_k]
=
\binom nk
2^{1-\binom{k}{2}}.
}
\]

If this value is below \(1\), a coloring with no monochromatic \(K_k\)
exists.

Ramsey theory extends to:

\[
\boxed{
\text{multiple colors},
\quad
\text{arbitrary graphs},
\quad
\text{hypergraphs},
\quad
\text{infinite sets},
\quad
\text{arithmetic patterns}.
}
\]

Examples include:

\[
\boxed{
\text{Erd\H{o}s--Szekeres},
\quad
\text{Van der Waerden},
\quad
\text{Schur}.
}
\]

The broad principle is

\[
\boxed{
\text{sufficient size}
\Longrightarrow
\text{unavoidable structured subconfiguration}.
}
\]

%%%%%%%%%%%%%%%%%%%%%%%%%%%%%%%%%%%%%%%%%%%%%%%%%%%%%%%%%%%%
\subsection{Section Connections}
\label{subsec:ramsey-theory-connections}
%%%%%%%%%%%%%%%%%%%%%%%%%%%%%%%%%%%%%%%%%%%%%%%%%%%%%%%%%%%%

\chapterconnection{
Ramsey Theory builds directly on Graph Theory and the pigeonhole
principle. It is closely connected with Extremal Combinatorics because
both study forbidden structures, while Probabilistic Combinatorics
provides many of the strongest general lower-bound methods. Algebraic
Combinatorics contributes explicit pseudorandom constructions and
spectral techniques. The next section, Design Theory, studies highly
regular finite incidence structures rather than unavoidable
substructures. Matroid Theory later abstracts independence and
dependence from graphs and vector spaces. Ramsey ideas also connect with
Number Theory through arithmetic progressions and additive equations,
with Geometry through Erd\H{o}s--Szekeres-type theorems, and with
Mathematical Logic through infinite Ramsey theory.
}

%%%%%%%%%%%%%%%%%%%%%%%%%%%%%%%%%%%%%%%%%%%%%%%%%%%%%%%%%%%%
% SECTION SOURCE TRACKING
%%%%%%%%%%%%%%%%%%%%%%%%%%%%%%%%%%%%%%%%%%%%%%%%%%%%%%%%%%%%

%------------------------------------------------------------
% 6.10 Ramsey Theory
%------------------------------------------------------------
%
% Source basis:
%   - Standard finite Ramsey theory.
%   - Standard graph Ramsey theory.
%   - Standard probabilistic Ramsey methods.
%   - Standard infinite and arithmetic Ramsey theory.
%
% Used for:
%   - basic Ramsey principle;
%   - party problem;
%   - two-color Ramsey numbers;
%   - diagonal and off-diagonal Ramsey numbers;
%   - symmetry of Ramsey numbers;
%   - trivial Ramsey values;
%   - proof of R(3,3)=6;
%   - Ramsey recurrence;
%   - refined parity recurrence;
%   - binomial upper bound;
%   - clique and independent-set formulation;
%   - probabilistic lower bounds;
%   - first-moment method;
%   - random-coloring arguments;
%   - graph Ramsey numbers;
%   - multicolor Ramsey numbers;
%   - infinite Ramsey theorem;
%   - hypergraph Ramsey theory;
%   - Erdős--Szekeres monotone subsequence theorem;
%   - geometric Ramsey phenomena;
%   - Van der Waerden's theorem;
%   - van der Waerden numbers;
%   - Schur's theorem;
%   - extremal and probabilistic connections;
%   - Ramsey graphs;
%   - pseudorandom constructions;
%   - Ramsey multiplicity;
%   - induced Ramsey theory;
%   - ordered Ramsey theory;
%   - canonical Ramsey theory;
%   - upper- and lower-bound proof strategies;
%   - applications and exercises.
%
% Notes:
%   - Probabilistic lower-bound methods build on Section 6.8.
%   - Graph terminology and clique/independent-set language build on
%     Section 6.9.
%   - Design Theory follows and shifts from unavoidable patterns to
%     highly regular incidence structures.
%   - Advanced arithmetic Ramsey theory connects strongly with Number
%     Theory and additive combinatorics.
%
%%%%%%%%%%%%%%%%%%%%%%%%%%%%%%%%%%%%%%%%%%%%%%%%%%%%%%%%%%%%